\section{Integral homology and Morse numbers}\label{morse:section}\label{sec:morse}

The construction requires two facts about Morse numbers: integral torsion
survives stabilization, and its resulting lower bound is attained by an
ordinary Morse function on the simply connected manifolds used here.
We prove the first fact directly for the precise condition at infinity
in the definition of $\SM$.  The second is the classical theorem of
Smale; a handle proof is included in Appendix~\ref{handle:appendix}.
The presentation follows the corresponding arguments in
\cite[Section~2 and Appendix~A]{OpenAICompanion}.

Recall that $t_i(W)$ is the minimum number of generators of
$\Tor H_i(W;\Z)$.  Put $r_i(W)=\operatorname{rank}H_i(W;\Z)$, so that
\[
 \lambda(W)=\sum_i r_i(W)+2\sum_i t_i(W).
\]
All homological indices are integers, and these quantities are zero
outside the homological range of $W$.

\begin{lemma}\label{morse:chain-bound}
Let $C_*$ be a finite chain complex of finitely generated free abelian
groups.  If $\rho_i$ is the rank of $H_i(C_*)$ and $\tau_i$ is the
minimum number of generators of its torsion subgroup, then
\[
 \operatorname{rank}C_i\geq\rho_i+\tau_i+\tau_{i-1},
 \qquad
 \sum_i\operatorname{rank}C_i\geq\sum_i\rho_i+2\sum_i\tau_i.
\]
\end{lemma}

\begin{proof}
Write $Z_i=\ker\partial_i$ and $B_i=\operatorname{im}\partial_{i+1}$.
These groups are free, and the two defining short exact sequences give
\begin{equation}\label{morse:chain-rank}
 \operatorname{rank}C_i
 =\rho_i+\operatorname{rank}B_i+\operatorname{rank}B_{i-1}.
\end{equation}
Smith normal form for $B_i\subset Z_i$ expresses
$Z_i/B_i=H_i(C_*)$ as a direct sum of a free group and cyclic groups,
one for each nonunit invariant factor of this inclusion.
Consequently $\operatorname{rank}B_i\geq\tau_i$.
Substitution in \eqref{morse:chain-rank} proves both inequalities.
\end{proof}

\begin{proposition}\label{morse:stable-lower}
For every connected closed smooth manifold $W$,
\[
 \SM(W)\geq\lambda(W).
\]
\end{proposition}

\begin{proof}
Let $F\colon W\times\R^k\to\R$ be an admissible Morse function.
We will construct a finite free integral chain complex with one
generator per critical point and with homology equal to that of $W$,
up to a degree shift.  For $k=0$, the ordinary Morse complex has these
properties.  Suppose $k>0$.

After a linear change of the auxiliary coordinates, the quadratic form
at infinity is
\[
 Q(u,v)=-|u|^2+|v|^2,
 \qquad (u,v)\in\R^d\times\R^{k-d}.
\]
Choose $R_0>0$ such that $F=Q$ for $|(u,v)|\geq R_0$ and every
critical point lies in $W\times B_{R_0}$.  Choose constants $A,T>0$ with
\[
 A>\max_{W\times\overline B_{R_0}}\{|F|,|Q|\},
 \qquad T>R_0,\qquad T^2>A.
\]
Set
\[
 K=W\times\overline B_T,\qquad
 K_-=K\cap\{F\leq-A\},\qquad
 K_+=K\cap\{F\leq A\}.
\]
The endpoint sublevels of $F$ and $Q$ agree:
\begin{equation}\label{morse:endpoint-pair}
 K_\pm=K\cap\{Q\leq\pm A\}.
\end{equation}
Indeed, outside the core the functions agree, while inside it both
sublevels at $-A$ are empty and both sublevels at $A$ contain the entire
core.

We next check that the truncation by $K$ introduces no extra cells.
On the side boundary $W\times\partial B_T$, the function is $Q$.
Its differential restricted to that boundary can vanish only at $u=0$
or $v=0$, with respective values $T^2$ and $-T^2$ when those loci
exist.  There are therefore no side-boundary critical points in the
level band $[-A,A]$.  A gradient-like vector field for $F$ on that
band can be chosen tangent to the side boundary: choose such a field
on the boundary using the nonvanishing restricted differential, extend
it to a collar, and combine it with an interior gradient-like field
by a partition of unity.  The condition $dF(V)>0$ away from critical
points is preserved by this combination.

The Morse attachment theorem now applies on this compact manifold with
corners.  Between critical values the tangent field gives the usual
deformation of sublevels, including at the side boundary; at each
critical point the change occurs in an interior Morse chart.
Thus, up to homotopy relative to $K_-$, the space $K_+$ is obtained
from $K_-$ by attaching one $j$-cell for each index-$j$ critical point
of $F$; see
\cite[Section~3]{Milnor1963} for the local attachment argument.
Coincident critical values can be separated by an arbitrarily small
perturbation supported in their Morse charts, without changing the
endpoint pair or the critical points and their indices.
Both endpoint sublevels have finite CW type.  Taking a finite CW model
for $K_-$ and cellular approximations of the attachments gives a finite
free relative cellular complex $C_*$ satisfying
\begin{equation}\label{morse:critical-complex}
 \sum_j\operatorname{rank}C_j=\#\Crit(F),
 \qquad H_j(C_*)\cong H_j(K_+,K_-;\Z).
\end{equation}

It remains to identify this homology.  The homotopy
$(u,v)\mapsto(u,(1-s)v)$, $0\leq s\leq1$, decreases both $Q$ and
the Euclidean norm.  By \eqref{morse:endpoint-pair}, it is a deformation
retraction of the pair onto
\[
 W\times\bigl(D_T^d,\{u\in D_T^d:|u|\geq\sqrt A\}\bigr).
\]
The disk--annulus pair has relative integral homology $\Z$ in degree
$d$ and zero in every other degree.  When $d=0$ it is the pair
$(\{0\},\varnothing)$, with the same assertion.  The relative
K\"unneth theorem therefore gives
\begin{equation}\label{morse:shifted-homology}
 H_j(C_*)\cong H_{j-d}(W;\Z).
\end{equation}
This calculation includes both definite signatures of $Q$.
Applying Lemma~\ref{morse:chain-bound} to
\eqref{morse:critical-complex}--\eqref{morse:shifted-homology} gives
$\#\Crit(F)\geq\lambda(W)$.  Taking the minimum over $F$ proves
the proposition.
\end{proof}

\begin{remark}\label{morse:infinity-conventions}
Some formulations of the stable Morse number allow
$\|dF-dQ\|$ to be uniformly bounded in the product of a metric on
$W$ and the Euclidean metric on $\R^k$.
This gives the same minimum as exact equality at infinity.
Indeed, write $h=F-Q$ and assume $\|dh\|\leq L$.
Compactness of $W$ gives $|h(x,z)|\leq K_0+L|z|$, while
nondegeneracy gives $\|dQ(z)\|\geq c|z|$ for some $c>0$.
Thus all critical points lie in $|z|\leq L/c$.
Choose a radial cutoff $\chi_R$ equal to one for $|z|\leq R$ and
zero for $|z|\geq2R$, with $\|d\chi_R\|\leq C/R$.
On the transition annulus,
\[
 \|d(\chi_Rh)\|\leq L+C(K_0/R+2L),
 \qquad \|dQ\|\geq cR.
\]
For sufficiently large $R$, the function $Q+\chi_Rh$ therefore
has exactly the original critical points, agrees with $F$ near them,
and equals $Q$ outside a compact set.  The converse inclusion is
immediate.  This also covers functions equal to $f(x)+Q(z)$ near
infinity, since $df$ is bounded on the closed manifold $W$.
\end{remark}

\begin{theorem}[Smale]\label{morse:equality}
Let $W$ be a simply connected closed smooth manifold of dimension
$D\geq9$.  There is an ordinary Morse function on $W$ with exactly
\[
 r_i(W)+t_i(W)+t_{i-1}(W)
\]
critical points of index $i$ for every $i$.  Consequently
\begin{equation}\label{morse:three-equal}
 \SM(W)=\Morse(W)=\lambda(W).
\end{equation}
\end{theorem}

\begin{proof}
The existence statement is the torsion-inclusive form of Smale's
Morse theorem \cite[Theorem~2.4 and the following paragraph,
pp.~134--135]{Smale1963}.  Appendix~\ref{handle:appendix} gives the
handle-cancellation proof in the stated dimension range, including the
case of handles of indices $2$ and $3$.
Summing the prescribed index counts gives $\lambda(W)$ critical points.
Since ordinary Morse functions are admissible in the definition of
$\SM$, Proposition~\ref{morse:stable-lower} yields
\[
 \lambda(W)\leq\SM(W)\leq\Morse(W)\leq\lambda(W).
\]
\end{proof}

The torsion count in these formulas combines different primes in the
same degree.  In particular, for distinct primes $p,q$,
\begin{equation}\label{morse:primary-maximum}
 d\bigl((\Z/p)^a\oplus(\Z/q)^b\bigr)=\max(a,b),
\end{equation}
where $d$ denotes the minimum number of generators.
Projection to each primary subgroup gives the lower bound, and pairing
cyclic summands by the Chinese remainder theorem gives the upper bound.
The degreewise maximum in \eqref{morse:primary-maximum} is the source of
the strict numerical comparison in our construction.
